\section{Conclusion}
\label{sec:conclusion}

This work asked how climate disruption changes urban accessibility and who loses the most from it. The answer rests on one operation. An environmental layer becomes a cost on the edges of a pedestrian network, and the framework that applies this cost is configured, not written for one city. Curitiba and Porto Alegre run the same operation with a temperature raster and two vector layers, each city naming its own layers and penalties. Supplying an environmental layer in the expected form is enough to run the analysis. The method transfers between cities, while the empirical result stays a property of each one.

The main finding is that how an edge fails matters as much as how many fail. A finite factor slows the network and keeps it whole, so the effect appears as longer travel times. An infinite factor removes the edge and breaks the network apart, so the effect appears as cells that can no longer reach a service. Both are settings of the same parameter, which makes the two effects comparable without changing the model. The socioeconomic result follows the same split. In Porto Alegre, the link with deprivation is strong and concentrated in sanitation. In Curitiba, it weakens once vegetation removes part of the heat and the weight shifts to income. Deprivation concentrates where exposed segments coincide with residents who have fewer alternatives. The two cities give different answers, and both are valid.

For planning, the practical reading follows from the same distinction. A corridor that slows needs shade, drainage or resurfacing, while a corridor that disappears needs a reinforced alternative route. Telling one from the other is the first step of an adaptation plan. The framework also returns the diagnosis at the scale where decisions are made, since a city average can stay low while the loss concentrates on the tracts least able to absorb it. Reading heat, water and slope through one interface puts them on a single accessibility metric, which lets a city rank its exposures instead of comparing unrelated risk scores. Every magnitude traces back to a stated source, so a planning office can see what rests on evidence and what is a choice.

Two steps would strengthen the work. The first is calibration against an observed event, since the indicators measure potential access over a static network and the demonstrative scenarios are not calibrated to observed magnitudes. The second is time variation, with exposures that rise, recede and leave the network changed. Both fit the declarative design as new parameters, and both move the framework from mapping vulnerability to supporting decisions. The contribution is the common ground for those steps, a way to ask which segments fail, by how much, and for whom, that works across cities and environmental layers.